\cotefile{104}
\scanpage{10}
la composition des $\operatorname{Hom}$ étant celle des $\mathrm{hom}$ \inserted{admiss.} entre les $ {\unsure{I|J}}_n$ (qui sont des immersions fermées). On a un foncteur canonique \deleted{\unsure{(défini|(définir}} \inserted{(\unsure{indépend.} de \unsure{tt} choix, \deleted{\unsure{à|s}} si on y tient)}
\[
M\longrightarrow M_0
\]
qui est une équivalence de catégories -- en fait, \deleted{\unsure{cel|et}} c'est un isom. de catégories, \unsure{NB|ssi} \deleted{\unsure{seulement}} $M$ est une catégorie réduite (deux objets \unsure{isom|isôm} sont identiques)

\deleted{\unsure{Je|La}} Je vais supposer pour la suite $M$ réduite. \unsure{on} peut \unsure{donc|en} l'identifier \unsure{:|à} la catégorie $M_0$, \unsure{définie|definie} par les données
\[
\resizebox{\linewidth}{!}{$\displaystyle
\left\{\begin{array}{@{}l@{}}
\begin{array}{@{}l@{\qquad}l@{\qquad}l@{}}
\mathbb N & (I_n)_{\unsure{n\in\mathbb N|\in\mathbb N}} & ,\;(\tau_n:I_n\to\mathbb N)_{n\in\mathbb N}\\
\text{ens ordonné} & \text{famille d'ens} & \text{famille d'\unsure{applic\textsuperscript{t}|applic\textsuperscript{n}}}\\
& \text{ordonnés, chaque }I_n\text{ un} & \text{strictement}\\
& \text{plus grand élément} & \text{croissantes}\\
&&\tau_n(I_n)=\mathbb N_{\leq n}\\
&&\tau_n^{-1}(n)=d_n\\
&&\text{(plus grd él. de }\unsure{?|I_n}
\end{array}\\[1ex]
(\unsure{\rho|\varphi}_{n,n'})_{\substack{(n',\unsure{n'|n})\in\mathbb N\times\mathbb N\\n'\leq n}}\\
\text{famille d'applications}\\
\begin{aligned}
\rho_{n,n'}:\tau_n^{-1}(n')&\longrightarrow\operatorname{Hom}(I_{n'},I_n)\\
\xi&\longmapsto (u^n_\xi:I_{n'}\unsure{-|\to}I_n)
\end{aligned}\\
\text{satisfaisant }\left\{\begin{aligned}
u^n_\xi &: I_{n'}\xrightarrow{\sim}(I_n)_{\leq\xi}&&
\text{(d'où }u^n_\xi(d_{n'})=\xi\text{)}\\
u^n_{\deleted{\unsure{\zeta|\eta}}\inserted{\xi}}\circ u^{n'}_\eta&=u^n_\zeta&&
\text{\unsure{avec|où} }\zeta=u^n_\xi(\eta)
\end{aligned}\right.\\[-.4ex]
\hfill\text{pour }\xi\in I_n,\ \eta\in I_{n'},\ n'=\tau_n(\xi),\ n''=\tau_{n'}(\eta)
\end{array}\right.
$}
\]
% archive sheet number: 9
\scanpage{11}
\mspage{3}
Inversement, de telles données définissent une catégorie $M_0$, satisfaisant $(i)(ii)$.

\uline{NB} Dans les cas que j'ai en tête, $\mathbb N=\unsure{\mathbb N|N}$

\uline{Ex 1} Catégorie des ens finis tot. ordonnés (ou : des \uline{simplexes} types $\Delta_n=\inserted{[0,n]}$) et des \unsure{applic\textsuperscript{tions}|applic\textsuperscript{ts}} strictement croissantes
% slanted
\margin{$ {\unsure{I|J}}_n={\unsure{\mathfrak P|\mathscr P}}^{*}(\Delta_n)$\\
$\operatorname{Card} I_n\unsure{-|=}2^{n+1}-1$}

\uline{Ex 2} Catégorie des \deleted{cubes} \inserted{``\uline{cubes}} épinglés'', \deleted{i.e.} dont \unsure{les|leurs} objets sont les \deleted{\unsure{systèmes|applications} \unsure{fléchés|fléchées}} cubes combinatoires,
\diagram{\begin{tikzcd}[ampersand replacement=\&, row sep=small,
execute at end picture={\draw[thin] (h104arrbot.south west) -- (h104arrtop.north east);}]
|[alias=h104arrtop]|\deleted{\widetilde A}\arrow[d] \\
|[alias=h104arrbot]|\deleted{A}
\end{tikzcd}}
\deleted{$A,\widetilde A$ ens finis non vides} \deleted{$\widetilde A\unsure{-|=}A$} épinglés par un ordre total sur l'ens des paires d'hyperfaces, \deleted{ou} et le choix d'une hyper\unsure{f.|face}, \unsure{ou face|dite} \unsure{notée|notion} $+$ dans chaque paire.
% slanted
\margin{$I_n={}$ens des sections partielles de $\unsure{(\!|\{}\,\Delta_{n-1}\times\{\pm1\}\to\Delta_{n-1}\unsure{)|\}}$ avec relation d'ordre \uline{opposée} de celle de prolongement\\
$\operatorname{Card} I_n=\deleted{\unsure{4}\,\unsure{2^{n}}\,\unsure{-2^{n}}\,\unsure{(2^{n+1}-1)}}3^n$}

\unsure{Tout|Un} épinglage du cube induit un épinglage de toute facette, donc \unsure{...} \deleted{\unsure{I|l}}

\uline{Ex. 3} Catégorie \deleted{des} \inserted{des} \uline{hémisphères}\deleted{\unsure{iques}} $D_n$, où $ {\unsure{D|I}}_n=(\Delta_n\times\{\pm1\})\amalg\{d_n\}$, avec $\{\pm1\}$ muni de l'ordre chaotique, $\deleted{\{}d_n\deleted{\}}$ \unsure{le|les} plus grand élément de $I_n$.
% slanted
\margin{$\operatorname{Card} I_n=2n+1$}

Dans \deleted{tous} ces \inserted{trois} exemples, on a

\textcircled{a} \inserted{Si $n\geq1$} La réalisation top. $|\deleted{\unsure{le}} {D}_n^*|$ de $I_n^*\inserted{=I_n\setminus\{d_n\}}$ est une \deleted{\unsure{antiq.|antique}}

\deleted{\textcircled{b} \unsure{La|les} \unsure{réalisation|réalisations}} $(n-1)$ sphère, (donc $|I_n|$ est un $n$-disque)
% archive sheet number: 10
\scanpage{12}
\textcircled{b} Les ens. \inserted{ordonnés} $ {\unsure{I|J}}_n^*$, (regardés comme $(n-1)$-poly\unsure{èdre|èdres}, sont \uline{réguliers}, i.e. le groupe des \unsure{autom} est transitif sur l'ens des \unsure{repères|épines} de $I_n^*$, i.e. l'ens des \unsure{chaines|chaînes} maximales
% slanted
\margin{b) est-il conséquence de a) ??}
\[
{\deleted{\unsure{x}}\unsure{\xi|x}}_0\unsure{<|\leftarrow}{\deleted{\unsure{x}}\xi}_1<\unsure{.\,-|\ldots}<{\deleted{\unsure{x}}\xi}_{n-1}
\]

\uline{Question} Y a-t-il d'autres exemples d'une suite \deleted{$(P_n\unsure{(\!-|=}I_{n+1}^{*})$} \deleted{$(\underset{\substack{\llcorner\mkern-5mu\overline{\phantom{I_n^*}}\mkern-5mu\lrcorner\\\inserted{P_{n-1}}}}{I_n^*})_{n\in\mathbb N,n\geq1}$} \inserted{$(P_n)_{n\geq1}$} d'ens ordonnés \unsure{présentant|possédant} les propriétés
% slanted
\margin{$P_n=I_{n+1}^{*}$}

a) $P_n$ est un \inserted{$n$-}polyèdre \uline{\unsure{régulier|réguliers}} sphérique

b) Si $\xi\in P_n$, \deleted{\unsure{est de dim. $n'$|est de dim., si}} \inserted{alors} $(P_n)_{\leq\xi}$ \deleted{est un polyèdre combinatoire de dimension $n'$, on a} $(P_n)_{\leq\xi}\simeq P_{n'}$, pour $n'$ convenable $\leq n$.

(\unsure{Ça|Sa} doit \unsure{être|etre} dans les livres de Coxeter !)

On aimerait de plus, surtout, être capables de trouver des $\deleted{\unsure{\rho_{\xi,n}|u^n_\xi}}\,{\unsure{\rho|\varphi}}_{n,n'}$ satisfaisant la condition de transitivité
% archive sheet number: 11
\scanpage{13}
\mspage{4}
Je considère maintenant la sous catégorie $ {\unsure{\mathscr B|\mathcal B|B}}_0$ de $\widehat M\supset M$, formée des objets $X$ de la forme suivante :
\[
\left(\begin{aligned}
&X\text{ est }\deleted{\text{\unsure{lim|limite}}}\inserted{\text{sup}}\text{ de ses sous-objets isomorphes \unsure{à|:} des}\\
&\deleted{\text{\unsure{ses|des}}}\qquad\text{objets de }M\text{ (modèles)}
\end{aligned}\right.
\]
% slanted
\margin{Cela \unsure{signifie|implique} que \unsure{tt} hom \deleted{de} $S\unsure{-|\to}X$, où $S\in\operatorname{Ob}M$, est un mono}

Si $K,L\subset X$ sont des sous-objets modèles, alors $K\cap L$ est sup des sous-objets de $K\cap L$ qui sont des modèles

Soit \unsure{donc|Donc} $\widetilde X$ l'ens des \deleted{\unsure{obj}} sous-objets de $X$ qui sont des modèles. Alors $\widetilde X$. est un ens ordonné \inserted{$K$}, muni de la structure suivante \inserted{(``$M$-épinglage'' de $K$)} : \deleted{\unsure{à|l'} \unsure{savoir|une}}

a) \deleted{\unsure{un|une} \unsure{de} \unsure{Pour}} Une application strictement croissante
\[
K\xrightarrow{\tau_K}\deleted{\unsure{I|N}}\unsure{N|\mathbb N}
\qquad\text{(application ``\uline{type}'')}
\]

b) Pour \unsure{tt} $\xi\in K$, \inserted{posant $n=\tau_K(\unsure{x|\xi})$,} \unsure{une|un} \deleted{\unsure{applic\textsuperscript{n}|applic\textsuperscript{t}}} \unsure{isomorphis\textsuperscript{me}|isomorphisme}
\[
\deleted{\unsure{D|I}}\;{\unsure{I|J}}_n\underset{u_\xi^K}{\simeq}K_{\leq\xi}
\qquad\text{(application d'épinglage de }K_{\leq\unsure{\xi|\eta}}\text{)}
\]
% slanted
\margin{les applications $I_n\unsure{-|\to}K$ définies par les $\xi\in K(n)=\unsure{(}\tau^{-1}(n)$ sont dite \uline{admissibles}}
\[
u_\xi^K(d_n)=\xi\quad\text{(= plus grd élément de }K_{\leq\xi}\text{)}
\]
avec \deleted{condition} les propriétés \deleted{de transitivité} :
% archive sheet number: 12
\scanpage{14}
\begingroup\makeatletter\tagsleft@true\makeatother
\[
\tau_K\circ u_\xi^K=\tau_n\tag*{(1)}
\]
et la propriété de transitivité pour les applications admissibles
\[
\begin{aligned}[t]
&\forall\,\xi\in K,\ \eta\in {\unsure{I|J}}_n\quad(n=\tau_K(\xi))\ \text{on a, si }n'=\tau_n(\eta)\\
&\qquad\deleted{\unsure{\tau|u}}\,u_\xi^K\circ u_\eta^n=u_\zeta^K,\quad\text{où }\zeta=u_\xi^K(\eta),
\end{aligned}\tag*{(2)}
\]
\endgroup
enfin

\deleted{(3) Si $\xi\in K,\ \eta\in K$, alors}

La catégorie dont les objets sont les objets admissibles de $\widehat M$, les flèches les \unsure{Isom|isom|mono} entre eux, est équivalente \unsure{:|à} celle des ens ordonnés $M$-épinglés.

D'autre part, les hom quelconques \inserted{dans $\unsure{\widehat C|\widehat M}$ entre tels objets} se réconstituent également, ce sont les applications \uline{strictement} \uline{croissantes}
\[
K\xrightarrow{f}L
\]
telles que $\left\{\begin{aligned}
&\text{a) }\tau_L\circ f=\tau_K\\
&\text{b) pour \unsure{tt} }\xi\in K,\ \inserted{\text{posant }\eta=f(\xi)\in L,}\ \text{on ait}
\end{aligned}\right.$ commutativité \unsure{du|de}
\diagram{\begin{tikzcd}[ampersand replacement=\&,column sep=small,row sep=large]
K_{\leq\xi}\arrow[rr,"{f|K_{\leq\xi}}"] \& \& {\deleted{K}L}_{\leq\eta}\\
\& I_n\arrow[ul,"\sim"']\arrow[ur,"\sim"] \& \& \text{où }n=\tau_K(\xi)=\tau_{\deleted{\unsure{K|\eta}}L}(\eta),
\end{tikzcd}}
% archive sheet number: 13
\scanpage{15}\mspage{5}
ce qui implique en particulier que $f\mid K_{\leq\unsure{\xi}}$ induit un iso
\[
K_{\leq\xi}\xrightarrow{\sim}{\deleted{\unsure{K|V}}L}_{\leq\eta}
\]

Nous nous intéressons surtout au cas où $f$ est injectif, ce qui correspond au \deleted{fait \unsure{que}} cas \unsure{d'un|d'une} \uline{mono} dans $\widehat C$.

Alors $f$ induit un \uline{isom}. de $K$ sur \deleted{\unsure{un} sous-ens} son image $L'=f(K)$, muni de l'ordre induit par $L$. Notons que $L'$ est une partie \uline{fermée} de $L$ $(x\in L',y\leq x\Rightarrow y\unsure{\in}L\unsure{,|\prime}$ et que \unsure{tte} partie fermée admet un $M$-épinglage \deleted{induit} unique, tel que l'inclusion soit admissible (épinglage induit) -- et \unsure{revenant} \unsure{:|à} $f$, bien sûr $K\overset{f}{\simeq}L'$ est compatible avec les $M$-épinglages.
% archive sheet number: 14
\scanpage{16}
Un cas intéressant est celui où $\deleted{\unsure{K|X}}\unsure{X|K}$ \deleted{\unsure{admet}} \inserted{est \unsure{le|la} sup \unsure{d'un|de un} \unsure{nb} fini de sous-}\deleted{un nb fini de \unsure{pts} maximaux, et}objets modèles, ce qui signifie que l'ens des éléments maximaux de $K\inserted{=\unsure{\widetilde X|X}}$ est fini, et que tt élément de $K$ est majoré par un élément maximal. Pour un ens. ordonné \deleted{ayant \unsure{ces}} \inserted{la $2^{\text{\unsure{ième|iem}}}$} \unsure{propriété|propriétés}, un $M$-épinglage est connu quand on connait les $u^K_{\deleted{\xi}\alpha}$ pour $\deleted{\xi}\alpha\inserted{\in K}$, maximal.

La condition pour que des donn.
\[
u^K_{\deleted{\xi}\alpha}:\ {\unsure{I|J}}_{\tau_{\unsure{n|K}}(\deleted{\xi}\alpha)}
   \xrightarrow{\sim} K_{\leq\deleted{\xi}\alpha}\qquad
   (\unsure{\xi}\in\operatorname{Max}K,
\]
définissent un \unsure{êpinglage|épinglage}, est que pour $\deleted{\unsure{\xi|\forall}}$ $\alpha,\beta\in\operatorname{Max}K$ et $\xi\in K$, $\xi\leq\alpha,\beta$, \deleted{posant} $n=\tau_K(\xi)$, \deleted{\unsure{des|les}} deux iso
\[
{\unsure{I|D}}_n\xrightarrow{\sim}K_{\leq\xi}
\]
définis par l'épinglage \deleted{\unsure{de|des}} $ {\deleted{\unsure{\tau}}u}^K_\alpha$ de $K_{\leq\alpha}$, et par l'épinglage $u^{\unsure{K|k}}_\beta$ de $K_{\leq\beta}$, \deleted{\unsure{soient}} les \unsure{mêmes}.
% archive sheet number: 15
\scanpage{17}\mspage{6}
La catégorie des objets admissibles de $\widehat C$ est stable par \deleted{\unsure{les} l'opération} sommes quelconques, et par sommes amalgamées $X\amalg_Z Y$, \unsure{don|dont} $Z\unsure{\mathord{-}|\to}X$ et $Z\mathord{-}Y$ des mono, et \unsure{le} foncteur $X\to\widetilde X$ commute \unsure{:|à} ces opérations, dans un sens évident. De plus, les objets de \unsure{type} fini donnent des objets de \unsure{t.f.}, pourvu, \unsure{(}dans le cas des sommes\unsure{)|,}, qu'on se borne \unsure{:|à} des sommes \uline{finies}.

\uline{NB} Quand les $\unsure{I_n|I}$ sont finis, alors les objets \inserted{adm.} de \unsure{t.f.} de $\widehat C$ \deleted{\unsure{aussi|sont} \unsure{t.f.|de} \unsure{les} $X$ \unsure{t.|,}} \inserted{correspondent aux} ens ordonnés $M$-\unsure{épinglés|êpinglés} $K$ qui \deleted{sont} \uline{finis}.
% archive sheet number: 16
